% Shared preamble of the paper. Section files must not define macros of their own.
\documentclass[11pt,a4paper]{amsart}

\usepackage[T1]{fontenc}
\usepackage[utf8]{inputenc}
\usepackage{lmodern}
\usepackage{amsmath,amssymb,amsthm,mathtools,mathrsfs}
\usepackage{enumitem}
\usepackage{booktabs,multicol}
\usepackage{xcolor}
\usepackage{tikz}
\usetikzlibrary{positioning,arrows.meta}
% Styles of the diagrams of the proof (Figures in the introduction and at the head of each part).
%   pfnode: a result proved in this paper     pfmain: one of Theorems A, B, C or the final theorem
%   pffact: a background result, admitted     pfpart: a block that stands for a whole part
%   pfedge: "is used in the proof of" (drawn from the premise to the conclusion)
\tikzset{
  pfbase/.style={draw, rounded corners=2pt, align=center, font=\scriptsize, inner sep=3pt,
                 text width=2.35cm, minimum height=0.9cm},
  pfnode/.style={pfbase},
  pfmain/.style={pfbase, line width=1pt},
  pffact/.style={pfbase, dashed, fill=black!8},
  pfpart/.style={pfbase, double, fill=black!3},
  pfedge/.style={-{Stealth[length=4pt]}, line width=0.4pt, draw=black!70},
}
\usepackage[margin=1.15in]{geometry}
\usepackage[colorlinks=true,linkcolor=blue,citecolor=blue,urlcolor=blue]{hyperref}

% Long Lean names leave few places to break a line; let TeX stretch instead of overflowing.
\emergencystretch=3em
\setcounter{tocdepth}{1}

% ---------------------------------------------------------------- theorem environments
\newtheorem{theorem}{Theorem}[section]
\newtheorem{proposition}[theorem]{Proposition}
\newtheorem{lemma}[theorem]{Lemma}
\newtheorem{corollary}[theorem]{Corollary}
% A result admitted without proof: it is proved in the Lean folder Lib and is known in the
% literature. State it precisely and give the source.
\newtheorem{fact}[theorem]{Fact}
\theoremstyle{definition}
\newtheorem{definition}[theorem]{Definition}
\newtheorem{notation}[theorem]{Notation}
\newtheorem{convention}[theorem]{Convention}
\newtheorem{example}[theorem]{Example}
\theoremstyle{remark}
\newtheorem{remark}[theorem]{Remark}
\numberwithin{equation}{section}

% ---------------------------------------------------------------- references to the Lean text
% \lean{Name}: a Lean declaration, by its full name.
% \leanfile{New/Reduction/MainInequality}: a Lean file, path relative to grh/src, no extension.
% \leanref{Name}: the same as \lean, flushed right, for the end of a statement.
%
% How a declaration is printed depends on whether an index has been loaded. The index is a file
% of lines \leanentry{Name}{number}{file}{line} (written by tools/make_leanindex.py). With it,
% \lean{Name} prints the short marker "L<number>", which links to the declaration on GitHub at
% the pinned commit, and \printleanindex prints the list of markers with the names. Without it
% (for instance when one section is compiled on its own) \lean{Name} prints the name in full.
\ExplSyntaxOn
\tl_new:N \l__lean_name_tl
\tl_const:Nx \c__lean_underscore_tl { \char_generate:nn { 95 } { 8 } }
\cs_generate_variant:Nn \tl_replace_all:Nnn { NVn }
\NewDocumentCommand{\leanfullname}{m}
 {
  \group_begin:
  \tl_set:Nn \l__lean_name_tl {#1}
  \tl_replace_all:NVn \l__lean_name_tl \c__lean_underscore_tl { \_ \allowbreak }
  \tl_replace_all:Nnn \l__lean_name_tl { . } { . \allowbreak }
  \tl_replace_all:Nnn \l__lean_name_tl { / } { / \allowbreak }
  \texttt { \small \tl_use:N \l__lean_name_tl }
  \group_end:
 }
\ExplSyntaxOff
\makeatletter
\newcommand{\leanrepo}{https://github.com/lsernau/grh-lean/blob/72fead102f663d5c7a407964a13dc0cf9d977207/grh/src}
\newcommand{\lean@list}{}
\newcommand{\leanentry}[4]{%
  \expandafter\gdef\csname lean@n@\detokenize{#1}\endcsname{#2}%
  \expandafter\gdef\csname lean@f@\detokenize{#1}\endcsname{#3}%
  \expandafter\gdef\csname lean@l@\detokenize{#1}\endcsname{#4}%
  \g@addto@macro\lean@list{\leanitem{#1}{#2}{#3}{#4}}}
\newcommand{\leanmark}[1]{\mbox{\normalfont\sffamily\small L#1}}
\NewDocumentCommand{\lean}{m}{%
  \@ifundefined{lean@n@\detokenize{#1}}%
    {\leanfullname{#1}}%
    {\begingroup\hypersetup{urlcolor=black!65}%
     \href{\leanrepo/\csname lean@f@\detokenize{#1}\endcsname.lean\#L\csname lean@l@\detokenize{#1}\endcsname}%
          {\leanmark{\csname lean@n@\detokenize{#1}\endcsname}}\endgroup}}
\newcommand{\leanfile}[1]{\leanfullname{#1}}
\newcommand{\leanref}[1]{\hfill\mbox{[\lean{#1}]}}
\newcommand{\leanitem}[4]{%
  \item[\leanmark{#2}\hfill] \href{\leanrepo/#3.lean\#L#4}{\leanfullname{#1}}}
\newcommand{\printleanindex}{%
  \begingroup\scriptsize
  \begin{multicols}{2}
  \begin{list}{}{\setlength{\leftmargin}{3.1em}\setlength{\labelwidth}{2.7em}\setlength{\labelsep}{0.4em}%
     \setlength{\itemsep}{0pt}\setlength{\parsep}{0pt}\setlength{\topsep}{0pt}\raggedright}
  \lean@list
  \end{list}
  \end{multicols}\endgroup}
\makeatother

% ---------------------------------------------------------------- Unicode in Lean names
% Lean names contain subscripts and Greek letters; make them printable with pdflatex.
\DeclareUnicodeCharacter{2080}{\textsubscript{0}}
\DeclareUnicodeCharacter{2081}{\textsubscript{1}}
\DeclareUnicodeCharacter{2082}{\textsubscript{2}}
\DeclareUnicodeCharacter{2083}{\textsubscript{3}}
\DeclareUnicodeCharacter{2084}{\textsubscript{4}}
\DeclareUnicodeCharacter{2085}{\textsubscript{5}}
\DeclareUnicodeCharacter{2086}{\textsubscript{6}}
\DeclareUnicodeCharacter{2087}{\textsubscript{7}}
\DeclareUnicodeCharacter{2088}{\textsubscript{8}}
\DeclareUnicodeCharacter{2089}{\textsubscript{9}}
\DeclareUnicodeCharacter{2090}{\textsubscript{a}}
\DeclareUnicodeCharacter{2091}{\textsubscript{e}}
\DeclareUnicodeCharacter{2092}{\textsubscript{o}}
\DeclareUnicodeCharacter{2093}{\textsubscript{x}}
\DeclareUnicodeCharacter{2095}{\textsubscript{h}}
\DeclareUnicodeCharacter{2096}{\textsubscript{k}}
\DeclareUnicodeCharacter{2097}{\textsubscript{l}}
\DeclareUnicodeCharacter{2098}{\textsubscript{m}}
\DeclareUnicodeCharacter{2099}{\textsubscript{n}}
\DeclareUnicodeCharacter{209A}{\textsubscript{p}}
\DeclareUnicodeCharacter{209B}{\textsubscript{s}}
\DeclareUnicodeCharacter{209C}{\textsubscript{t}}
\DeclareUnicodeCharacter{1D62}{\textsubscript{i}}
\DeclareUnicodeCharacter{2C7C}{\textsubscript{j}}
\DeclareUnicodeCharacter{207A}{\textsuperscript{+}}
\DeclareUnicodeCharacter{207B}{\textsuperscript{-}}
\DeclareUnicodeCharacter{2070}{\textsuperscript{0}}
\DeclareUnicodeCharacter{2074}{\textsuperscript{4}}
\DeclareUnicodeCharacter{207F}{\textsuperscript{n}}
\DeclareUnicodeCharacter{03B1}{\ensuremath{\alpha}}
\DeclareUnicodeCharacter{03B2}{\ensuremath{\beta}}
\DeclareUnicodeCharacter{03B3}{\ensuremath{\gamma}}
\DeclareUnicodeCharacter{03B4}{\ensuremath{\delta}}
\DeclareUnicodeCharacter{03B5}{\ensuremath{\varepsilon}}
\DeclareUnicodeCharacter{03B6}{\ensuremath{\zeta}}
\DeclareUnicodeCharacter{03B7}{\ensuremath{\eta}}
\DeclareUnicodeCharacter{03B8}{\ensuremath{\theta}}
\DeclareUnicodeCharacter{03D1}{\ensuremath{\vartheta}}
\DeclareUnicodeCharacter{03B9}{\ensuremath{\iota}}
\DeclareUnicodeCharacter{03BA}{\ensuremath{\kappa}}
\DeclareUnicodeCharacter{03BB}{\ensuremath{\lambda}}
\DeclareUnicodeCharacter{03BC}{\ensuremath{\mu}}
\DeclareUnicodeCharacter{03BD}{\ensuremath{\nu}}
\DeclareUnicodeCharacter{03BE}{\ensuremath{\xi}}
\DeclareUnicodeCharacter{03C0}{\ensuremath{\pi}}
\DeclareUnicodeCharacter{03C1}{\ensuremath{\rho}}
\DeclareUnicodeCharacter{03F1}{\ensuremath{\varrho}}
\DeclareUnicodeCharacter{03C3}{\ensuremath{\sigma}}
\DeclareUnicodeCharacter{03C4}{\ensuremath{\tau}}
\DeclareUnicodeCharacter{03C5}{\ensuremath{\upsilon}}
\DeclareUnicodeCharacter{03C6}{\ensuremath{\varphi}}
\DeclareUnicodeCharacter{03D5}{\ensuremath{\phi}}
\DeclareUnicodeCharacter{03C7}{\ensuremath{\chi}}
\DeclareUnicodeCharacter{03C8}{\ensuremath{\psi}}
\DeclareUnicodeCharacter{03C9}{\ensuremath{\omega}}
\DeclareUnicodeCharacter{0393}{\ensuremath{\Gamma}}
\DeclareUnicodeCharacter{0394}{\ensuremath{\Delta}}
\DeclareUnicodeCharacter{0398}{\ensuremath{\Theta}}
\DeclareUnicodeCharacter{039B}{\ensuremath{\Lambda}}
\DeclareUnicodeCharacter{039E}{\ensuremath{\Xi}}
\DeclareUnicodeCharacter{03A0}{\ensuremath{\Pi}}
\DeclareUnicodeCharacter{03A3}{\ensuremath{\Sigma}}
\DeclareUnicodeCharacter{03A6}{\ensuremath{\Phi}}
\DeclareUnicodeCharacter{03A8}{\ensuremath{\Psi}}
\DeclareUnicodeCharacter{03A9}{\ensuremath{\Omega}}
\DeclareUnicodeCharacter{2113}{\ensuremath{\ell}}
\DeclareUnicodeCharacter{2102}{\ensuremath{\mathbb{C}}}
\DeclareUnicodeCharacter{211D}{\ensuremath{\mathbb{R}}}
\DeclareUnicodeCharacter{211A}{\ensuremath{\mathbb{Q}}}
\DeclareUnicodeCharacter{2124}{\ensuremath{\mathbb{Z}}}
\DeclareUnicodeCharacter{2115}{\ensuremath{\mathbb{N}}}
\DeclareUnicodeCharacter{1D538}{\ensuremath{\mathbb{A}}}
\DeclareUnicodeCharacter{1D53D}{\ensuremath{\mathbb{F}}}
\DeclareUnicodeCharacter{1D4AE}{\ensuremath{\mathcal{S}}}
\DeclareUnicodeCharacter{1D4AA}{\ensuremath{\mathcal{O}}}
\DeclareUnicodeCharacter{1D4AB}{\ensuremath{\mathcal{P}}}
\DeclareUnicodeCharacter{1D52D}{\ensuremath{\mathfrak{p}}}
\DeclareUnicodeCharacter{27E8}{\ensuremath{\langle}}
\DeclareUnicodeCharacter{27E9}{\ensuremath{\rangle}}
\DeclareUnicodeCharacter{266D}{\ensuremath{\flat}}
\DeclareUnicodeCharacter{2192}{\ensuremath{\to}}
\DeclareUnicodeCharacter{21A6}{\ensuremath{\mapsto}}
\DeclareUnicodeCharacter{2264}{\ensuremath{\le}}
\DeclareUnicodeCharacter{2265}{\ensuremath{\ge}}
\DeclareUnicodeCharacter{2260}{\ensuremath{\ne}}
\DeclareUnicodeCharacter{2208}{\ensuremath{\in}}
\DeclareUnicodeCharacter{2211}{\ensuremath{\sum}}
\DeclareUnicodeCharacter{220F}{\ensuremath{\prod}}
\DeclareUnicodeCharacter{2282}{\ensuremath{\subset}}
\DeclareUnicodeCharacter{2286}{\ensuremath{\subseteq}}
\DeclareUnicodeCharacter{221E}{\ensuremath{\infty}}
\DeclareUnicodeCharacter{2261}{\ensuremath{\equiv}}
\DeclareUnicodeCharacter{2223}{\ensuremath{\mid}}
\DeclareUnicodeCharacter{2016}{\ensuremath{\|}}
\DeclareUnicodeCharacter{230A}{\ensuremath{\lfloor}}
\DeclareUnicodeCharacter{230B}{\ensuremath{\rfloor}}
\DeclareUnicodeCharacter{2308}{\ensuremath{\lceil}}
\DeclareUnicodeCharacter{2309}{\ensuremath{\rceil}}
\DeclareUnicodeCharacter{221A}{\ensuremath{\surd}}
\DeclareUnicodeCharacter{2218}{\ensuremath{\circ}}
\DeclareUnicodeCharacter{2297}{\ensuremath{\otimes}}
\DeclareUnicodeCharacter{2295}{\ensuremath{\oplus}}
\DeclareUnicodeCharacter{2200}{\ensuremath{\forall}}
\DeclareUnicodeCharacter{2203}{\ensuremath{\exists}}
\DeclareUnicodeCharacter{2227}{\ensuremath{\wedge}}
\DeclareUnicodeCharacter{2228}{\ensuremath{\vee}}
\DeclareUnicodeCharacter{00AC}{\ensuremath{\neg}}
\DeclareUnicodeCharacter{2243}{\ensuremath{\simeq}}
\DeclareUnicodeCharacter{2245}{\ensuremath{\cong}}
\DeclareUnicodeCharacter{2229}{\ensuremath{\cap}}
\DeclareUnicodeCharacter{222A}{\ensuremath{\cup}}
\DeclareUnicodeCharacter{22A4}{\ensuremath{\top}}
\DeclareUnicodeCharacter{22A5}{\ensuremath{\bot}}
\DeclareUnicodeCharacter{2202}{\ensuremath{\partial}}
\DeclareUnicodeCharacter{2032}{\ensuremath{{}'}}

\DeclareUnicodeCharacter{2212}{\ensuremath{-}}
\DeclareUnicodeCharacter{2075}{\textsuperscript{5}}
\DeclareUnicodeCharacter{2076}{\textsuperscript{6}}
\DeclareUnicodeCharacter{2077}{\textsuperscript{7}}
\DeclareUnicodeCharacter{2078}{\textsuperscript{8}}
\DeclareUnicodeCharacter{2079}{\textsuperscript{9}}
\DeclareUnicodeCharacter{1EF8}{\ensuremath{\tilde{Y}}}
\DeclareUnicodeCharacter{0168}{\ensuremath{\tilde{U}}}
\DeclareUnicodeCharacter{0304}{\ensuremath{\bar{\ }}}
\DeclareUnicodeCharacter{1FB1}{\ensuremath{\bar\alpha}}
\DeclareUnicodeCharacter{03C2}{\ensuremath{\varsigma}}
\DeclareUnicodeCharacter{1D52F}{\ensuremath{\mathfrak{r}}}
\DeclareUnicodeCharacter{2262}{\ensuremath{\not\equiv}}
\DeclareUnicodeCharacter{266E}{\ensuremath{\natural}}
\DeclareUnicodeCharacter{1D5A5}{\ensuremath{\mathsf{F}}}
\DeclareUnicodeCharacter{2205}{\ensuremath{\emptyset}}

\DeclareUnicodeCharacter{266F}{\ensuremath{\sharp}}
\DeclareUnicodeCharacter{208A}{\textsubscript{+}}
\DeclareUnicodeCharacter{1D51E}{\ensuremath{\mathfrak{a}}}
\DeclareUnicodeCharacter{210B}{\ensuremath{\mathcal{H}}}
\DeclareUnicodeCharacter{0303}{\ensuremath{\tilde{\ }}}

% ---------------------------------------------------------------- number systems
\newcommand{\C}{\mathbb{C}}
\newcommand{\R}{\mathbb{R}}
\newcommand{\Q}{\mathbb{Q}}
\newcommand{\Z}{\mathbb{Z}}
\newcommand{\N}{\mathbb{N}}
\newcommand{\F}{\mathbb{F}}
\newcommand{\A}{\mathbb{A}}
\newcommand{\Af}{\mathbb{A}_{\mathrm{f}}}
\newcommand{\Qbar}{\overline{\mathbb{Q}}}
\newcommand{\Ql}{\mathbb{Q}_\ell}
\newcommand{\Qlbar}{\overline{\mathbb{Q}}_\ell}
\newcommand{\Zl}{\mathbb{Z}_\ell}
\newcommand{\Qp}{\mathbb{Q}_p}
\newcommand{\Zp}{\mathbb{Z}_p}
\newcommand{\Cp}{\mathbb{C}_p}
\newcommand{\Cl}{\mathbb{C}_\ell}
\newcommand{\cO}{\mathcal{O}}
\newcommand{\fp}{\mathfrak{p}}
\newcommand{\fq}{\mathfrak{q}}
\newcommand{\fa}{\mathfrak{a}}
\newcommand{\fm}{\mathfrak{m}}

% ---------------------------------------------------------------- operators
\renewcommand{\Re}{\operatorname{Re}}
\renewcommand{\Im}{\operatorname{Im}}
\DeclareMathOperator{\GL}{GL}
\DeclareMathOperator{\SL}{SL}
\DeclareMathOperator{\U}{U}
\DeclareMathOperator{\Sp}{Sp}
\DeclareMathOperator{\Gal}{Gal}
\DeclareMathOperator{\Hom}{Hom}
\DeclareMathOperator{\End}{End}
\DeclareMathOperator{\Aut}{Aut}
\DeclareMathOperator{\tr}{tr}
\DeclareMathOperator{\Tr}{Tr}
\DeclareMathOperator{\Nm}{Nm}
\DeclareMathOperator{\rk}{rk}
\DeclareMathOperator{\ord}{ord}
\DeclareMathOperator{\vol}{vol}
\DeclareMathOperator{\covol}{covol}
\DeclareMathOperator{\disc}{disc}
\DeclareMathOperator{\Frob}{Frob}
\DeclareMathOperator{\id}{id}
\DeclareMathOperator{\Ind}{Ind}
\DeclareMathOperator{\Res}{Res}
\DeclareMathOperator{\Sym}{Sym}
\DeclareMathOperator{\Span}{span}
\DeclareMathOperator{\image}{im}
\DeclareMathOperator{\coker}{coker}
\DeclareMathOperator{\sgn}{sgn}
\DeclareMathOperator{\diag}{diag}
\DeclareMathOperator{\length}{length}
\DeclareMathOperator{\supp}{supp}
\DeclareMathOperator{\Stab}{Stab}
\DeclareMathOperator{\Lie}{Lie}
\DeclareMathOperator{\Spec}{Spec}
\DeclareMathOperator{\Ext}{Ext}
\DeclareMathOperator{\height}{ht}

\DeclarePairedDelimiter{\abs}{\lvert}{\rvert}
\DeclarePairedDelimiter{\norm}{\lVert}{\rVert}
\DeclarePairedDelimiter{\floor}{\lfloor}{\rfloor}
\DeclarePairedDelimiter{\ceil}{\lceil}{\rceil}
\DeclarePairedDelimiter{\pair}{\langle}{\rangle}
\newcommand{\set}[1]{\{\, #1 \,\}}
\newcommand{\eps}{\varepsilon}
\newcommand{\isom}{\xrightarrow{\ \sim\ }}
\newcommand{\inj}{\hookrightarrow}
\newcommand{\surj}{\twoheadrightarrow}

% ---------------------------------------------------------------- letters of this proof
% (see the notation table in Section 3)
\newcommand{\chiK}{\chi_K}
\newcommand{\LambdaK}{\Lambda_K}
\newcommand{\sinf}{s_\infty}
\newcommand{\sfin}{s_{\mathrm{fin}}}
\newcommand{\smain}{s_{\mathrm{main}}}
\newcommand{\hmain}{h_{\mathrm{main}}}
\newcommand{\hdir}{h_{\mathrm{dir}}}
\newcommand{\rem}{\mathrm{rem}}
\newcommand{\rdE}{\mathrm{rd}_E}
\newcommand{\dT}{d_T}
\newcommand{\Gbar}{\overline{G}}
\newcommand{\Poin}{\mathscr{P}}
\newcommand{\nPoin}{\mathfrak{p}}
